\section{Butterworth Filter Design and Sensitivity}

\subsection{Butterworth Prototypes}

\subsubsection{Butterworth magnitude response}
\[
|H(j\omega)|^2=\frac{1}{1+\left(\frac{\omega}{\omega_c}\right)^{2N}}.
\]
A Butterworth filter is maximally flat in the passband. At \(\omega=\omega_c\), the magnitude is \(1/\sqrt{2}\), or \(-3\) dB.

\subsubsection{Butterworth pole angles}
\[
p_k=\omega_c e^{j\left(\frac{\pi}{2}+\frac{(2k-1)\pi}{2N}\right)},\qquad k=1,\ldots,N.
\]
The stable Butterworth poles lie on a circle of radius \(\omega_c\) in the left half-plane.

\subsubsection{Butterworth second-order stage}
\[
H_k(s)=\frac{\omega_c^2}{s^2+a_k\omega_c s+\omega_c^2},
\qquad
a_k=2\zeta_k.
\]
Even-order Butterworth filters are cascades of second-order sections. The stage coefficient \(a_k\) sets damping and quality factor.

\subsubsection{Common Butterworth coefficients}
\[
\begin{array}{c|c}
N & \text{second-order }a_k\text{ values}\\ \hline
2 & 1.4142\\
4 & 0.7654,\ 1.8478\\
6 & 0.5176,\ 1.4142,\ 1.9319\\
8 & 0.3902,\ 1.1111,\ 1.6629,\ 1.9616
\end{array}
\]
These coefficients are used when decomposing normalized Butterworth filters into cascaded second-order stages.

\subsection{Sallen-Key Lowpass}

\subsubsection{Sallen-Key lowpass transfer function}
\[
H_{\mathrm{SK,LP}}(s)=
\frac{K}{s^2R_1R_2C_1C_2+s\left(R_1C_1+R_1C_2+R_2C_2(1-K)\right)+1}.
\]
The Sallen-Key lowpass realizes a second-order denominator using two resistors, two capacitors, and an amplifier gain \(K\).

\subsubsection{Sallen-Key natural frequency and damping coefficient}
\[
\omega_n=\frac{1}{\sqrt{R_1R_2C_1C_2}},
\]
\[
a=2\zeta=\frac{R_1C_1+R_1C_2+R_2C_2(1-K)}{\sqrt{R_1R_2C_1C_2}}.
\]
Component products set the natural frequency. Component ratios and gain set damping.

\subsubsection{Unity-gain Sallen-Key lowpass}
\[
K=1 \quad \Rightarrow \quad
a=\frac{R_1C_1+R_1C_2}{\sqrt{R_1R_2C_1C_2}}.
\]
Unity-gain design removes the \(R_2C_2(1-K)\) term and is used in the low-sensitivity design slides.

\subsection{Sensitivity}

\subsubsection{Relative sensitivity}
\[
S_x^y=\frac{\partial y/y}{\partial x/x}=\frac{x}{y}\frac{\partial y}{\partial x}.
\]
Sensitivity measures fractional output change per fractional component change.

\subsubsection{Sensitivity of products and powers}
\[
y=Kx_1^{\alpha_1}x_2^{\alpha_2}
\quad \Rightarrow \quad
S_{x_i}^{y}=\alpha_i.
\]
Power-law expressions give sensitivities directly from exponents.

\subsubsection{Monte Carlo component tolerance}
\[
x_i=x_{i,\mathrm{nom}}(1+\epsilon_i).
\]
Monte Carlo analysis perturbs component values within their tolerance ranges to estimate spread in cutoff frequency, gain, and damping.

\subsection{Scaling}

\subsubsection{Frequency scaling of RC filters}
\[
\omega_n' = K_F\omega_n,\qquad
R'C'=\frac{RC}{K_F}.
\]
Increasing cutoff frequency requires reducing the relevant \(RC\) products.

\subsubsection{Impedance scaling}
\[
R'=K_ZR,\qquad C'=\frac{C}{K_Z}.
\]
Impedance scaling changes practical component values while preserving the same \(RC\) products and transfer function.
